\documentclass[11pt]{article}

\usepackage{amsmath,amssymb}
\usepackage{xurl}
\usepackage[hidelinks]{hyperref}
\setlength{\emergencystretch}{2em}

% Shared commands for notes in docs/paper-gaps/.
%
% Two halves.  The link commands below are project identity: OWNER/REPO is the
% placeholder that scripts/bootstrap_project.py replaces with this project's
% GitHub slug and Pages site.  Everything after them is NOTATION, and notation
% belongs to the source paper: the definitions here are an example set, and a
% project replaces them with the macros of its own paper's style file.  The one
% rule that never changes: a command defined here must mean exactly what the
% source paper's macro means, or a note silently says something else.

\newcommand{\Lean}{\textsc{Lean}}
\newcommand{\leanid}[1]{\textsf{#1}}
\newcommand{\ghissue}[1]{\href{https://github.com/OWNER/REPO/issues/#1}{GitHub issue \##1}}
\newcommand{\ghpr}[1]{\href{https://github.com/OWNER/REPO/pull/#1}{PR \##1}}
% Local issue tracker (issues/NNNN-*.md in this repository); no remote URL.
\newcommand{\localissue}[1]{issue \##1 (\path{issues/})}
\newcommand{\doclink}[1]{\href{https://OWNER.github.io/REPO/docs/#1.html}{\path{#1}}}

% --- Notation: replace with the source paper's own macros --------------------
% \F, \N, \C, \R, \tr, \ot, \eps, \E, \bx, \bu, \bv, \by

\newcommand{\F}{\mathbb{F}}
\newcommand{\N}{\mathbb{N}}
\newcommand{\C}{\mathbb{C}}
\newcommand{\R}{\mathbb{R}}
\newcommand{\tr}{\mathrm{tr}}
\newcommand{\eps}{\epsilon}
\newcommand{\ot}{\otimes}
\newcommand{\E}{\mathop{\bf E\/}}

% bold random variables
\newcommand{\bu}{\boldsymbol{u}}
\newcommand{\bv}{\boldsymbol{v}}
\newcommand{\bx}{{\boldsymbol{x}}}
\newcommand{\by}{\boldsymbol{y}}

% T1 so literal < and > in \texttt placeholders typeset as themselves
\usepackage[T1]{fontenc}

% bra-ket
\usepackage{braket}

% --- Example structured spaces ---
\newcommand{\polyfunc}[3]{\mathcal{P}(#1, #2, #3)}
\newcommand{\polymeas}[3]{\mathrm{PolyMeas}(#1, #2, #3)}
\newcommand{\polysub}[3]{\mathrm{PolySub}(#1, #2, #3)}

% Approximation relations (paper uses \approx_\eps and \simeq_\zeta)
\newcommand{\approxeps}{\approx_{\varepsilon}}
\newcommand{\simgeq}{\simeq_{\zeta}}

\renewcommand{\ghpr}[1]{%
  \href{https://github.com/Dengnifer/MIPStarRE-A/pull/#1}{PR \##1}}

\title{A Field-Valued Construction of the Combined Point Measurements}
\author{}
\date{2026-09-09}

\begin{document}
\maketitle

\section{At a glance}

We give a replacement proof of the combined-point lemma in the Pauli basis
analysis of \cite{Ji2020MIPStar}.  The source combines binary refinements and
then invokes quantum linearity; we instead combine the field-valued point
measurements directly.  Parseval's identity prevents any loss depending on
the field size, and orthonormalization produces projective measurements on
the existing expanded local spaces.  The resulting squared-distance error is
$K\eps^{1/8}$ for a universal constant $K$.  No additional ancilla is needed.

\paragraph{Difficulty and estimated weight.}
The construction is proved.  Its formalization is of medium difficulty and
occupies roughly two thousand lines.

\paragraph{Mathlib/project split and key mathematical inputs.}
The proof is predominantly project-specific placement, Fourier, and
measurement estimates.  It reuses the project's rounding and distance
theorems, with Mathlib supplying finite sums, inner products, matrix
adjoints, and real powers.  This note documents a change of proof, not a
change of the source theorem.\footnote{The implementation is in
\path{MIPStarRE/QPBT/Combining/Points.lean} and its supporting modules under
\path{MIPStarRE/QPBT/Combining/Points/}.  See
\ghpr{410},
review findings F2--F3.}

\section{Key theorem forms}

Let $q=2^t$ and let the admissible Pauli basis test parameters and a projective
strategy of success probability at least $1-\eps$ be fixed.  Write
$E=(\C^q)^{\otimes M}$ for the ancillary register space and
$\ket{\Phi_E}$ for its normalized EPR vector.  We work with the expanded state
\[
  \ket{\hat\psi}
  =\ket{\psi}_{AB}\otimes\ket{\Phi_E}_{A'A''}
    \otimes\ket{\Phi_E}_{B'B''}.
\]
The original player spaces $H_A$ and $H_B$ may differ.  For an opposite pair
of placements $(AA',BA'')$ or $(AB'',BB')$, in either order, use subscripts
$1,2$ to indicate the placements.  For a fixed point pair $(x,z)$, abbreviate
the expanded point projections by $X_i(a)=\hat M^{(\mathrm{Point},X),x}_a$
and $Z_i(b)=\hat M^{(\mathrm{Point},Z),z}_b$ at placement $i$.
Operators on opposite placements commute; we do not assume that $X_i(a)$
commutes with $Z_i(b)$.

For families indexed by $(x,z,a,b)$ define
\[
 d_{\hat\psi}^2(A,B)
 =\E_{x,z}\sum_{a,b\in\F_q}
       \|(A^{x,z}_{a,b}-B^{x,z}_{a,b})\ket{\hat\psi}\|^2,
\]
where the points are uniform and independent.  The source asserts the
existence of projective measurements $Q^{x,z}$ on each $H_i\otimes E$ with
polynomially small self-consistency error and closeness to both ordered
products.\footnote{Lemma \path{lem:qld-4-10} and displays
\path{eq:qld-q-self-cons}, \path{eq:qld-q-cons-m-hat-xz}, and
\path{eq:qld-q-cons-m-hat-zx} in
\path{references/qpbt-paper/14_analysis_of_the_pauli_basis_test.tex},
lines 689--709.}
The replacement proof gives, for every directed opposite pair,
\begin{equation}\label{eq:target}
 \max\{d_{\hat\psi}^2(Q_1,Q_2),\,
        d_{\hat\psi}^2(Q_1,X_2Z_2),\,
        d_{\hat\psi}^2(Q_1,Z_2X_2)\}
 \leq K\eps^{1/8}.
\end{equation}
Here products retain the common outcome pair $(a,b)$, and $K$ is independent
of $q$, the number of points, and the Hilbert-space dimensions.

\section{The source construction and the replacement}

The source first constructs a sandwich POVM for binary trace refinements,
rounds it, proves approximate linearity in the trace parameters $(r,s)$,
and applies the quantum linearity theorem separately for each $(x,z)$.
This last application returns an extended space.  The source absorbs it by
assuming sufficiently many zero-state ancillas.  To formalize that argument
on a common extended space, one must make the ancillary spaces uniform over
the point pairs and identify the padded state.\footnote{The three-step
construction begins at line 711; the binary sandwich and its rounding are
at lines 731--785, and the linearity application with zero-state padding is
at lines 825--832 of the same paper mirror.  The remaining common-extension
construction for that route is discussed in
\path{docs/paper-gaps/qpbt_linearity-theorem-quotation.tex}.}
The replacement below does not apply the linearity theorem.  We therefore
need neither its ancillary spaces nor a common-extension construction.

\paragraph{Parseval transfer.}
Put $\chi(u)=(-1)^{\operatorname{tr}(u)}$, where
$\operatorname{tr}:\F_q\to\F_2$ is the field trace.  The expanded observables
$\mathcal X_i^r(x)$ and $\mathcal Z_i^s(z)$ give the Fourier formulas
$X_i(a)=q^{-1}\sum_r\chi(ar)\mathcal X_i^r(x)$ and
$Z_i(b)=q^{-1}\sum_s\chi(bs)\mathcal Z_i^s(z)$.
Character orthogonality on $\F_q^2$, applied to vectors in the state space,
gives the exact identity
\begin{equation}\label{eq:parseval}
 \sum_{a,b}\|[X_i(a),Z_i(b)]\ket{\hat\psi}\|^2
 =\frac{1}{q^2}\sum_{r,s}
    \|[\mathcal X_i^r(x),\mathcal Z_i^s(z)]\ket{\hat\psi}\|^2.
\end{equation}
Indeed, expansion of the squared norm yields character sums equal to
$q^2$ for equal frequency pairs and zero otherwise, cancelling two of the
four inverse powers of $q$.  The right-hand side is thus an average, not
an unnormalized sum.  The observable commutation bound underlying the
source's binary-refinement estimate implies that the average of either side
is $O(\sqrt\eps)$.\footnote{See \path{lem:qld-comm-cons} and its proof,
paper lines 466--505.  The exact identity is
\leanid{sum\_norm\_place\_expPointOp\_commutator\_sq}; its averaged consequence
is \leanid{expPoint\_comm}, in \path{Points/Commutation.lean}.}

\paragraph{The field-valued sandwich.}
Set $R_i(a,b)=Z_i(b)X_i(a)Z_i(b)$.  These operators are positive because
$R_i(a,b)=(X_i(a)Z_i(b))^*(X_i(a)Z_i(b))$, and they sum to the identity.
Moreover,
$R_i(a,b)-Z_i(b)X_i(a)=Z_i(b)[X_i(a),Z_i(b)]$.
Projection contraction and \eqref{eq:parseval} therefore give
$d_{\hat\psi}^2(R_i,Z_iX_i)=O(\sqrt\eps)$.

For a fixed $(x,z)$, define $D^X_a=X_1(a)-X_2(a)$,
$D^Z_b=Z_1(b)-Z_2(b)$, $W_b=Z_1(b)Z_2(b)$, and
$K^i_{a,b}=[X_i(a),Z_i(b)]$.
Two applications of the projective overlap identity give
\[
 1-\sum_{a,b}\langle\hat\psi,R_1(a,b)R_2(a,b)\hat\psi\rangle
 =\frac12\sum_b\|D^Z_b\hat\psi\|^2
  +\frac12\sum_{a,b}\|D^X_a W_b\hat\psi\|^2.
\]
The left-hand side is the off-diagonal consistency mass.  To bound the second
sum, use only opposite-placement commutation to write
\[
 D^X_aW_b=W_bD^X_a+Z_2(b)K^1_{a,b}-Z_1(b)K^2_{a,b}.
\]
The squared norm of three terms is at most three times the sum of their
squared norms.  Projection contraction and
$\sum_b\|Z_2(b)v\|^2=\|v\|^2$ then bound the consistency mass by
\[
 c_{x,z}=\frac12\sum_b\|D^Z_b\hat\psi\|^2
  +\frac32\left(\sum_a\|D^X_a\hat\psi\|^2
    +\sum_{a,b}\|K^1_{a,b}\hat\psi\|^2
    +\sum_{a,b}\|K^2_{a,b}\hat\psi\|^2\right).
\]
Point self-consistency bounds the first two sums on average by $O(\eps)$;
Parseval bounds the commutator sums by $O(\sqrt\eps)$.  Consequently
$\E_{x,z}c_{x,z}\leq C_0(\eps+\sqrt\eps)$ with a universal $C_0$.
This exact overlap calculation replaces the source's Cauchy--Schwarz
chain.\footnote{Displays \path{eq:qld-rw-self-cons-1} through
\path{eq:qld-rw-self-cons-4}, paper lines 743--773.  The formal identities
are \leanid{sandwich\_overlap\_identity} and
\leanid{sandwich\_defect\_pointwise\_le}; the bound is
\leanid{sandwichDefectBound}, in \path{Points/Consistency.lean}.}

\paragraph{Rounding and the error exponent.}
For two POVMs on opposite factors of a normalized bipartite state, the
orthonormalization lemma gives a projective measurement on the first
\emph{existing} factor at squared distance at most $220c^{1/4}$ from its
POVM when their consistency defect is at most $c\geq0$.  Apply this lemma
for each point pair, first with $AA'$ as the local factor and then with
$BA''$ as the local factor, placing the unused EPR pair on the other side.
Both constructions return measurements on $H_A\otimes E$ and
$H_B\otimes E$, respectively.  Jensen's inequality gives
\[
 d_{\hat\psi}^2(Q_i,R_i)
 \leq220\,\E_{x,z}c_{x,z}^{1/4}
 \leq220\bigl(\E_{x,z}c_{x,z}\bigr)^{1/4}
 \leq220\bigl(C_0(\eps+\sqrt\eps)\bigr)^{1/4}.
\]
This step uses the same orthonormalization result as the source, now with
outcomes in $\F_q^2$.  In Lean the averaging inequality follows by applying
the square-root Cauchy--Schwarz estimate twice and identifying
$c^{1/4}=\sqrt{\sqrt c}$.\footnote{The rounding input is
\leanid{projective\_rounding\_with\_explicit\_constant} in
\path{MIPStarRE/QPBT/Games/DistanceTheorems/ProjectiveRounding.lean}.
Its use and the averaging lemma \leanid{avgOver\_rpow\_quarter\_le} are in
\path{Points/Orthonormalization.lean}.}

Transporting the two point factors one at a time gives
$d_{\hat\psi}^2(Z_1X_1,X_2Z_2)=O(\eps)$ by their self-consistency and
projective completeness.  Chaining this estimate with rounding, sandwich
closeness, and commutation gives all three distances in \eqref{eq:target}
bounded by a universal multiple of
$(\eps+\sqrt\eps)^{1/4}+\sqrt\eps+\eps$.
For $0\leq\eps\leq1$ this is $O(\eps^{1/8})$.  For $\eps>1$, each distance
is at most $4$: both projective families and both ordered-product families
satisfy $\sum_{a,b}A_{a,b}^*A_{a,b}\leq I$.
Enlarging $K$ therefore proves \eqref{eq:target} for every admissible error.

\section{Placements and formal comparison}

The auxiliary permutation $\tau=(A'\ B'')(A''\ B')$ fixes $A,B$ and
exchanges the two EPR pairs while reversing each pair.  Hence its unitary
fixes $\ket{\hat\psi}$ for an arbitrary strategy, and exchanges $AA'$ with
$AB''$ and $BA''$ with $BB'$.  It transfers the rounding bounds to the
second bipartition scheme without exchanging the players or assuming a
symmetric strategy.  We prove the estimates in both directions on opposite
placements using the corresponding point consistency estimates.

This involution is distinct from the product of the blueprint permutations
$\sigma=(A\ B)(A'\ B')(A''\ B'')$ and
$\theta=(A\ B)(A'\ A'')$: their product acts as a four-cycle on the
ancillas.  The auxiliary permutation supports symmetric equivalents but
does not identify with that product or prove the blueprint's full transfer
lemma.\footnote{See \path{def:symmetric-equivalents} and
\path{lem:symmetric-equivalents-transfer} in
\path{blueprint/src/chapter/ch14_qpbt_observables.tex}.
The auxiliary permutation is \leanid{eprCrossSwap} in
\path{Points/Placement.lean}.}

The combined-point statement in the blueprint and Lean retains the original
local spaces, both product orders, and all directed opposite comparisons.
The blueprint presents the source's binary-refinement proof; the formal
construction uses the replacement above.\footnote{See
\path{lem:qld-4-10} in
\path{blueprint/src/chapter/ch15_qpbt_combining.tex} and
\leanid{exists\_combinedPointsWitness} in
\path{MIPStarRE/QPBT/Combining/Points.lean}.}
No hypothesis supplying a common ancilla, approximate linearity, or exact
same-player commutation is added.  Finite field models and finite nonempty
register types encode the source's ambient spaces.

\section{Conclusion}

There is no theorem-statement discrepancy in this construction.  The
field-valued proof realizes the source's polynomial-error conclusion with a
dimension-independent bound $K\eps^{1/8}$ on the original expanded spaces.
It bypasses, but does not discharge, the common-extension obligation of the
separate quantum-linearity route.  The crosswise EPR exchange is an auxiliary
involution, not the product of permutations named in the blueprint.

\bibliographystyle{alpha}
\bibliography{references}

\end{document}
